% ============================================================================
%  Modelo UTXO frente a modelo de cuentas
%  Módulo 08 · Modelos de estado — ChainNotes
%
%  Versión LaTeX del apunte web temas/utxo-vs-cuentas.html
%  Los números de selección de monedas salen del mismo modelo de tarifas que
%  usa assets/js/utxo-vs-cuentas.js (entrada 68 vB, salida 31 vB, LTR 10).
%
%  Compilar:  pdflatex utxo-vs-cuentas.tex   (dos pasadas, por el índice)
% ============================================================================

\documentclass[11pt,a4paper]{article}

% --- idioma y codificación -------------------------------------------------
\usepackage[T1]{fontenc}
\usepackage[utf8]{inputenc}
\usepackage[spanish,es-noshorthands,es-tabla]{babel}
\usepackage{lmodern}
\usepackage{microtype}

% --- página ----------------------------------------------------------------
\usepackage[a4paper,margin=2.6cm,headheight=14pt]{geometry}
\usepackage{fancyhdr}
\usepackage{enumitem}
\usepackage{booktabs}
\usepackage{tabularx}
\usepackage{array}
\usepackage{caption}
\usepackage{longtable}

% --- matemáticas -----------------------------------------------------------
\usepackage{amsmath,amssymb,amsthm}
\usepackage{mathtools}

% --- gráficas --------------------------------------------------------------
\usepackage{tikz}
\usetikzlibrary{positioning,arrows.meta,calc,fit,backgrounds,shapes.geometric,
                decorations.pathreplacing}

% --- enlaces y cajas (hyperref antes que tcolorbox) ------------------------
\usepackage[colorlinks=true,linkcolor=violetdeep,urlcolor=violetdeep,
            citecolor=violetdeep,pdftitle={Modelo UTXO frente a modelo de cuentas},
            pdfauthor={ChainNotes}]{hyperref}
\usepackage[most]{tcolorbox}

% ---------------------------------------------------------------------------
%  Paleta (heredada de assets/css/styles.css)
% ---------------------------------------------------------------------------
\definecolor{violetc}{HTML}{7A68F5}
\definecolor{violetdeep}{HTML}{5B46E0}
\definecolor{mintc}{HTML}{1FA88A}
\definecolor{coralc}{HTML}{E04E5C}
\definecolor{amberc}{HTML}{C98A1F}
\definecolor{skyc}{HTML}{1E88C7}
\definecolor{inkc}{HTML}{282344}

\hypersetup{linkcolor=violetdeep}

% ---------------------------------------------------------------------------
%  Estilo de párrafo y encabezados
% ---------------------------------------------------------------------------
\setlength{\parindent}{0pt}
\setlength{\parskip}{0.55em}
\captionsetup{font=small,labelfont={bf,color=violetdeep}}

\renewcommand{\topfraction}{0.9}
\renewcommand{\bottomfraction}{0.8}
\renewcommand{\textfraction}{0.07}
\renewcommand{\floatpagefraction}{0.75}
\setcounter{topnumber}{3}
\setcounter{totalnumber}{4}

\pagestyle{fancy}
\fancyhf{}
\fancyhead[L]{\small\color{inkc}\textsc{UTXO frente a cuentas}}
\fancyhead[R]{\small\color{inkc}\thepage}
\renewcommand{\headrulewidth}{0.4pt}

% ---------------------------------------------------------------------------
%  Entornos de teoremas
% ---------------------------------------------------------------------------
\theoremstyle{definition}
\newtheorem{definicion}{Definición}
\theoremstyle{plain}
\newtheorem{proposicion}{Proposición}
\theoremstyle{remark}
\newtheorem*{observacion}{Observación}

% ---------------------------------------------------------------------------
%  Cajas de color
% ---------------------------------------------------------------------------
\newtcolorbox{nota}[1]{enhanced,breakable,colback=skyc!6,colframe=skyc!75!black,
  coltitle=white,fonttitle=\bfseries\small,title={#1},
  left=7pt,right=7pt,top=5pt,bottom=5pt,boxrule=0.7pt,arc=2.5pt,
  before upper={\setlength{\parskip}{0.5em}}}

\newtcolorbox{caso}[1]{enhanced,breakable,colback=amberc!8,colframe=amberc!85!black,
  coltitle=white,fonttitle=\bfseries\small,title={#1},
  left=7pt,right=7pt,top=5pt,bottom=5pt,boxrule=0.7pt,arc=2.5pt,
  before upper={\setlength{\parskip}{0.5em}}}

\newtcolorbox{aviso}[1]{enhanced,breakable,colback=coralc!7,colframe=coralc!85!black,
  coltitle=white,fonttitle=\bfseries\small,title={#1},
  left=7pt,right=7pt,top=5pt,bottom=5pt,boxrule=0.7pt,arc=2.5pt,
  before upper={\setlength{\parskip}{0.5em}}}

\newtcolorbox{clave}[1]{enhanced,breakable,colback=mintc!7,colframe=mintc!85!black,
  coltitle=white,fonttitle=\bfseries\small,title={#1},
  left=7pt,right=7pt,top=5pt,bottom=5pt,boxrule=0.7pt,arc=2.5pt,
  before upper={\setlength{\parskip}{0.5em}}}

\newtcolorbox{laboratorio}[1]{enhanced,breakable,colback=violetc!6,colframe=violetc!80!black,
  coltitle=white,fonttitle=\bfseries\small,title={#1},
  left=7pt,right=7pt,top=5pt,bottom=5pt,boxrule=0.7pt,arc=2.5pt,
  before upper={\setlength{\parskip}{0.5em}}}

\newtcolorbox{algoritmo}[1]{enhanced,breakable,colback=inkc!3,colframe=inkc!60,
  coltitle=white,fonttitle=\bfseries\small,title={#1},
  left=9pt,right=7pt,top=5pt,bottom=5pt,boxrule=0.6pt,arc=2.5pt}

% ---------------------------------------------------------------------------
%  Atajos de notación
% ---------------------------------------------------------------------------
\newcommand{\si}{\textcolor{mintc}{\checkmark}}
\newcommand{\no}{\textcolor{coralc}{$\times$}}
\newcommand{\mecanismo}[1]{\textbf{\textcolor{violetdeep}{#1}}}

\setlist[itemize]{label=\textcolor{violetc}{$\bullet$}}

\newlist{opciones}{enumerate}{1}
\setlist[opciones]{label=(\alph*),leftmargin=2.4em,itemsep=0.1em,topsep=0.2em}

\begin{document}

\makeatletter
\DeclareRobustCommand{\%}{\char`\%\relax}
\makeatother

\begin{titlepage}
\centering
\vspace*{2.5cm}

{\Large\color{violetdeep}\textsc{ChainNotes}\par}
{\small\color{inkc}Módulo 08 · Modelos de estado\par}

\vspace{1.4cm}
{\Huge\bfseries Modelo UTXO\par}
\vspace{0.4cm}
{\LARGE frente a modelo de cuentas\par}

\vspace{1.2cm}
\begin{minipage}{0.78\textwidth}
\centering\itshape\color{inkc}
Las funciones hash dicen quién autoriza y el sellado de tiempo dice cuándo. Falta
una pregunta anterior y más básica: cuando un nodo afirma que tienes 8 unidades,
¿qué hay escrito realmente en su disco? Hay dos respuestas razonables, y cada
gran familia de blockchains escogió una.
\end{minipage}

\vspace{1.4cm}
\begin{tikzpicture}[scale=0.95]
  % objetos
  \node[inkc,font=\small\bfseries] at (-3.1,1.5) {objetos};
  \draw[violetc,line width=0.8pt,fill=violetc!15,rounded corners=3pt] (-4.1,0.55) rectangle (-3.3,1.15);
  \node[violetdeep,font=\scriptsize] at (-3.7,0.85) {5};
  \draw[violetc,line width=0.8pt,fill=violetc!15,rounded corners=3pt] (-2.9,0.55) rectangle (-2.1,1.15);
  \node[violetdeep,font=\scriptsize] at (-2.5,0.85) {3};
  \node[inkc,font=\scriptsize] at (-3.1,0.1) {dos monedas};
  % celda
  \node[inkc,font=\small\bfseries] at (3.1,1.5) {celda};
  \draw[skyc,line width=0.8pt,fill=skyc!12,rounded corners=3pt] (2.0,0.55) rectangle (4.2,1.15);
  \node[skyc,font=\scriptsize] at (3.1,0.85) {0xa1c3\ldots\ $\mapsto$ 8};
  \node[inkc,font=\scriptsize] at (3.1,0.1) {un renglón};
  \node[violetdeep,font=\small\itshape] at (0,0.85) {«tienes 8»};
\end{tikzpicture}

\vfill
{\small Apuntes de Blockchain · Proyecto académico\par}
{\small\color{inkc}\today\par}
\end{titlepage}

\tableofcontents
\clearpage

\begin{nota}{Cómo usar este documento}
Es la versión imprimible del apunte interactivo
\texttt{temas/utxo-vs-cuentas.html}. Los tres laboratorios aparecen aquí como
ejemplos resueltos con números fijos, calculados con el mismo modelo de tarifas
que usa \texttt{assets/js/utxo-vs-cuentas.js}. Al final hay seis ejercicios con
soluciones en el apéndice~\ref{app:sol}.
\end{nota}

% ===========================================================================
\section{La pregunta que los separa}\label{sec:pregunta}
% ===========================================================================

Un nodo no almacena una idea abstracta de riqueza: almacena bytes. Y para los
bytes que representan «Alicia tiene 8» hay dos formas sensatas de organizarlos.

\textbf{Como un conjunto de objetos.} El nodo guarda una colección de monedas
indivisibles, cada una con un valor y una condición que dice quién puede
gastarla. Que Alicia tenga 8 significa que existen por ahí un objeto de 5 y otro
de 3 que solo ella puede consumir. No hay ningún renglón con su nombre. Y los
objetos jamás se editan: se destruyen y nacen otros.

\textbf{Como una tabla de saldos.} El nodo guarda un diccionario que lleva de
dirección a saldo. Que Alicia tenga 8 significa que en la celda de su dirección
está escrito un 8. No hay monedas, hay una celda. Y las celdas jamás se
destruyen: se sobrescriben.

\begin{clave}{Efectivo y banca, con su letra chica}
UTXO se parece al efectivo: pagas 6 entregando un billete de 5 y uno de 3, y
recibes 2 de cambio; no existe medio billete. Cuentas se parece al banco: el
cajero tacha tu saldo y escribe otro, sin que ningún objeto físico cambie de
manos.

La analogía tiene fecha de caducidad: en el efectivo real las denominaciones son
fijas, mientras que un UTXO puede valer cualquier cantidad. El cambio que recibes
no es un billete estándar, es uno fabricado a la medida en ese mismo instante.
\end{clave}

La diferencia de fondo se ve en cómo se obtiene el saldo. En el modelo de objetos
no está escrito en ninguna parte: es una consulta.
\[
  \mathrm{saldo}(\text{alicia}) = \sum \{\, v : (v,c) \in U,\ c \text{ la satisface alicia} \,\}
\]
En el modelo de celdas es el dato, y leerlo cuesta una búsqueda:
\[
  \mathrm{saldo}(\text{alicia}) = S[\text{alicia}].\text{saldo}
\]
Derivar el saldo o leerlo: esa es la raíz de todo lo que sigue.

% ===========================================================================
\section{Definiciones}\label{sec:def}
% ===========================================================================

\begin{definicion}[UTXO]
\emph{Unspent Transaction Output}, salida de transacción no gastada. Es una
pareja $(\text{valor}, \text{condición de gasto})$ identificada por el par
$(\textit{txid}, \text{índice})$, llamado \textbf{outpoint}.

El estado global es el conjunto $U$ de todos los UTXO vigentes. Una transacción
es una función que toma un subconjunto finito de $U$, lo elimina e inserta
salidas nuevas:
\[
  U' = (U \setminus \text{entradas}) \cup \text{salidas}.
\]
\textbf{Validez:} cada entrada existe en $U$, su condición de gasto se satisface,
y $\sum \text{entradas} \ge \sum \text{salidas}$. La diferencia es la comisión.
\end{definicion}

\begin{definicion}[Cuentas]
El estado global es una función parcial
$S: \text{dirección} \to (\text{nonce}, \text{saldo}, \text{código}, \text{almacenamiento})$.
Ethereum distingue las cuentas controladas por una llave privada de las cuentas
de contrato, que además tienen código y almacenamiento propio.

Una transacción modifica entradas de $S$:
\[
  S' = S[\text{origen} \mapsto -\text{monto}-\text{gas},\ \text{destino} \mapsto +\text{monto}].
\]
\textbf{Validez:} la firma corresponde al origen, el saldo alcanza, y el nonce de
la transacción coincide exactamente con el de la cuenta.
\end{definicion}

\begin{observacion}[La asimetría que hay que notar]
En UTXO la validez de una transacción se decide mirando \emph{solo} las entradas
que consume: la verificación es local. En cuentas hay que mirar el renglón
completo del emisor, incluido un contador que no tiene nada que ver con el
dinero. Ese contador —el nonce— hace un trabajo que en UTXO nadie necesita hacer.
\end{observacion}

% ===========================================================================
\section{Anatomía de una transacción}\label{sec:anatomia}
% ===========================================================================

\begin{algoritmo}{Bitcoin}
\begin{verbatim}
version    2
entradas
  [0] 3f9c...a1:0   testigo: <firma> <pubkey>
  [1] 77b2...e4:1   testigo: <firma> <pubkey>
salidas
  [0] 600 000 sat -> OP_1 <beto>
  [1] 190 000 sat -> OP_1 <alicia_cambio>
locktime   0
\end{verbatim}
No aparece ningún remitente: el emisor se deduce de quién pudo firmar las
entradas. Tampoco aparece la comisión, que es el residuo aritmético.
\end{algoritmo}

\begin{algoritmo}{Ethereum}
\begin{verbatim}
chainId    1
nonce      7
to         0x5aE2...c19f
value      6 000 000 000 000 000 000
gasLimit   21 000
maxFeePerGas          24 gwei
maxPriorityFeePerGas   1 gwei
data       0x
firma      v, r, s
\end{verbatim}
Todo es explícito: destinatario, monto, límite de gas, cadena. Y un nonce, el
único campo cuya razón de ser no se adivina a primera vista.
\end{algoritmo}

\begin{aviso}{La consecuencia inmediata}
En Bitcoin el tamaño de la transacción depende de cuántas monedas hayas tenido
que juntar, y la comisión depende del tamaño. \textbf{Pagar la misma cantidad
puede costarte el doble si tu cartera está fragmentada.} En Ethereum una
transferencia simple cuesta 21\,000 de gas sin importar tu historial.
\end{aviso}

% ===========================================================================
\section{Elegir monedas es NP-difícil}\label{sec:seleccion}
% ===========================================================================

Tu cartera tiene varias monedas y quieres pagar una cantidad concreta. ¿Cuáles
entregas? Cada moneda que incluyas agranda la transacción y encarece la comisión.
Pero si usas pocas monedas grandes, el cambio que recibes fragmenta tu cartera
para el futuro. Es un problema con dos horizontes temporales en conflicto.

\begin{nota}{Modelo de tarifas de este documento}
\begin{center}
\small
\begin{tabular}{@{}lr@{}}
\toprule
Entrada & 68 vB \\
Salida & 31 vB \\
Encabezado & 11 vB \\
Tasa de largo plazo & 10 sat/vB \\
Umbral de polvo & 546 sat \\
\bottomrule
\end{tabular}
\end{center}
Son los pesos de una entrada y una salida modernas, redondeados.
\end{nota}

\subsection{La métrica de desperdicio}

Bitcoin Core no elige un algoritmo: los ejecuta todos y compara. Desde la versión
22.0 (2021) cada candidato se puntúa con una métrica llamada \emph{waste}:
\[
  \text{desperdicio} = \text{peso}_{\text{entradas}} \times
  (\text{tasa} - \text{tasa}_{\text{LP}}) + \text{costo}_{\text{cambio}} + \text{exceso}.
\]

El término que manda es el primero. Si la tasa actual está por debajo de la de
largo plazo, gastar monedas de más es barato ahora y caro después, así que la
métrica empuja a \textbf{consolidar}. Si la tasa está por arriba, empuja a usar
el \textbf{mínimo de entradas}. La cartera se comporta distinto según el mercado
de comisiones sin que el usuario toque nada.

El \textbf{valor efectivo} de una moneda es lo que aporta ya descontado lo que
cuesta gastarla:
\[
  \text{valor efectivo} = \text{valor} - 68\,\text{vB} \times \text{tasa}.
\]
Una moneda con valor efectivo negativo \emph{cuesta dinero} incluirla. A
20 sat/vB eso ocurre por debajo de 1360 sat: el polvo deja de ser gastable sin
que nadie lo borre.

\subsection{Los algoritmos}

\textbf{Branch and Bound.} Propuesto por Murch en 2016 e incorporado en Bitcoin
Core 0.17. Busca un subconjunto cuyo valor efectivo caiga dentro de una ventana
estrecha alrededor del objetivo, de modo que no haga falta crear cambio. Una
transacción sin cambio ahorra los 31 vB de la salida ahora y los 68 de gastarla
después, corta el grafo de transacciones y mejora la privacidad.

\textbf{CoinGrinder.} Añadido en Bitcoin Core 27.0 (abril de 2024). Busca el
conjunto de entradas de peso mínimo y sí produce cambio. Solo se activa cuando la
tasa supera unas tres veces la de largo plazo. La restricción es deliberada: si
minimizaras entradas todo el tiempo, triturarías la cartera. En simulaciones
sobre los picos de 2023 redujo el gasto entre 4\% y 9.5\%.

\textbf{Mochila y sorteo aleatorio.} El \emph{knapsack} heredado hace mil
intentos aleatorios. El sorteo aleatorio simple (SRD) toma monedas al azar: es
peor en comisión pero excelente en privacidad, porque no deja huella reconocible.

\begin{laboratorio}{Ejemplo resuelto}
Cartera de 1, 3, 7, 12 y 20 mil satoshis. Se quiere pagar 15 mil a 20 sat/vB. El
costo del cambio es $31 \times 20 + 68 \times 10 = 1300$ sat. Hay 20 subconjuntos
válidos; los tres mejores:

\begin{center}
\small
\begin{tabular}{@{}lrrlr@{}}
\toprule
Selección & Entradas & Comisión & Cambio & Desperdicio \\
\midrule
$7000 + 12\,000$          & 2 & 4000 & sin cambio & \textbf{1800} \\
$20\,000$                 & 1 & 2820 & 2180       & 1980 \\
$1000 + 7000 + 12\,000$   & 3 & 5000 & sin cambio & 2120 \\
\bottomrule
\end{tabular}
\end{center}

Gana $\{7000, 12\,000\}$ sin cambio. El cálculo:
$2 \times 68 \times (20-10) = 1360$ del término temporal, más 440 de exceso que
se va a comisión.
\end{laboratorio}

\begin{observacion}
La opción de una sola entrada paga menos comisión \emph{hoy} —2820 frente a
4000— y aun así pierde, porque arrastra el costo del cambio: 1300 sat que habrá
que pagar mañana. Es justo el punto de la métrica.
\end{observacion}

\begin{aviso}{Por qué NP-difícil}
Decidir si existe un subconjunto de monedas cuya suma sea exactamente el objetivo
es el problema de la \textbf{suma de subconjuntos}. Lo que salva a las carteras
reales es que el número de monedas es pequeño y que basta una solución
suficientemente buena, no la óptima.
\end{aviso}

% ===========================================================================
\section{El nonce hace el trabajo que en UTXO hace el objeto}\label{sec:nonce}
% ===========================================================================

Aquí está la simetría profunda entre los dos modelos.

En \textbf{UTXO}, una transacción firmada es imposible de repetir: consume
objetos concretos que después de la primera vez ya no existen. Reenviarla no hace
nada, porque sus entradas no están en $U$. La protección es automática y no
requiere ningún campo adicional.

En \textbf{cuentas}, una transferencia firmada de 6 unidades es un mensaje válido
para siempre: los saldos son números, no objetos, y nada impide aplicarla cien
veces. Por eso cada cuenta lleva un contador que se incrementa con cada
transacción, y cada transacción declara el valor que espera encontrar.

\begin{aviso}{Lo que cuesta el contador}
El nonce impone un \textbf{orden total estricto} a todo lo que hace una cuenta.
Si la transacción 5 se queda atorada por comisión baja, la 6 y la 7 no pueden
ejecutarse aunque paguen bien. En UTXO no existe tal cosa: dos pagos que gastan
monedas distintas son independientes aunque salgan de la misma cartera.
\end{aviso}

\begin{observacion}[Un matiz histórico]
El nonce protege contra la repetición \emph{dentro} de una cadena, pero no
\emph{entre} cadenas. Cuando Ethereum se bifurcó en 2016, la misma transacción
firmada era válida en las dos redes. La solución fue el \textbf{EIP-155}, que
metió el identificador de cadena dentro de lo que se firma.
\end{observacion}

% ===========================================================================
\section{Paralelismo}\label{sec:paralelismo}
% ===========================================================================

Dos transacciones se pueden validar en paralelo si no tocan el mismo estado. En
UTXO, «el mismo estado» significa la misma moneda: dos pagos al mismo
destinatario son independientes, porque cada uno crea su propia moneda nueva. En
el modelo de cuentas, los dos escriben el mismo renglón, y hay conflicto.

\begin{laboratorio}{Seis pagos, dos modelos}
Los pagos: P1 Ana$\to$Zoe, P2 Beto$\to$Zoe, P3 Caro$\to$Zoe, P4 Dan$\to$Eva,
P5 Eva$\to$Ana, P6 Zoe$\to$Beto.

\textbf{Como UTXO:} un solo carril. Los seis gastan monedas distintas, así que se
validan a la vez. Que P1, P2 y P3 paguen los tres a Zoe da igual, porque cada uno
crea una moneda nueva e independiente.

\textbf{Como cuentas:} cuatro carriles. Zoe aparece en P1, P2, P3 y P6, y los
cuatro escriben su renglón, así que hay que serializarlos.
\end{laboratorio}

El número de carriles no depende de cuánto trabajo hay, sino de cómo está
repartido el estado. Los seis pagos cuestan lo mismo en los dos modelos; lo que
cambia es cuántos caben a la vez.

\begin{observacion}[La letra chica]
El ejemplo está construido para que el contraste sea nítido. En la práctica
muchas transacciones de Bitcoin gastan salidas creadas por otras transacciones
del mismo bloque, lo que sí crea dependencias. Pero la asimetría estructural es
real, y es la que llevó a varios diseños recientes a recuperar parte del modelo
de objetos: Sui trata cada activo como un objeto con dueño explícito; Aptos y
Monad conservan cuentas pero ejecutan de forma optimista con Block-STM; Solana
exige que cada transacción declare por adelantado qué estado va a tocar.
\end{observacion}

% ===========================================================================
\section{El estado en disco}\label{sec:estado}
% ===========================================================================

\subsection{El conjunto UTXO}

Bitcoin Core lo mantiene en una base clave-valor local, el \textbf{chainstate},
indexada por outpoint. Es el único estado que un nodo necesita para validar
bloques nuevos.

Su tamaño es una preocupación viva: pasó de unos 52 millones de salidas a finales
de 2017 a alrededor de 173 millones a principios de 2025. Nadie puede borrar
entradas salvo gastándolas, así que el polvo acumulado es permanente.

Demostrar que una transacción está en un bloque es una rama de Merkle corta.
Demostrar que una salida \emph{sigue sin gastarse} requiere estructuras
adicionales, y es un problema abierto en el que se trabaja con acumuladores como
\textbf{Utreexo}.

\subsection{El árbol de estado}

Ethereum usa un \textbf{árbol de Merkle Patricia hexario}: cada nivel ramifica
según un nibble de la clave. La raíz va en cada encabezado de bloque, lo que
permite demostrar criptográficamente el saldo de cualquier cuenta.

El costo es el tamaño de las pruebas: una prueba de acceso ronda los 3 kB, y con
miles de accesos por bloque el testigo total se vuelve inviable de difundir en un
intervalo de doce segundos.

\begin{nota}{El giro de 2025}
Durante años la solución candidata fueron los \textbf{árboles Verkle}. Ganaban en
tamaño de prueba, pero sus compromisos polinomiales descansan en curvas
elípticas, e introducirlos en el corazón del protocolo significaba apoyarse en
criptografía que no resiste a un adversario cuántico, justo cuando NIST
recomienda abandonarla hacia 2030.

La propuesta vigente es el \textbf{EIP-7864}, un árbol binario unificado basado
solo en funciones hash, con ramas unas cuatro veces más cortas. Se sacrificó
tamaño de prueba a cambio de una estructura que probablemente no habrá que volver
a cambiar.
\end{nota}

% ===========================================================================
\section{Privacidad}\label{sec:privacidad}
% ===========================================================================

En el modelo de cuentas tu dirección es un identificador permanente: toda tu
historia cuelga de un solo renglón público. En UTXO no hay renglón, y las
carteras generan una dirección nueva para cada cambio, así que a primera vista el
rastro se disuelve.

A primera vista. La \textbf{heurística de propiedad común de entradas} dice que
si una transacción gasta varias monedas a la vez, lo más probable es que todas
pertenezcan al mismo dueño, porque solo él pudo firmarlas todas. Basta esa
observación para volver a unir buena parte de las monedas de una cartera.

\begin{clave}{La selección de monedas vuelve a aparecer}
Elegir monedas deja de ser solo un problema de comisiones y se vuelve también uno
de privacidad. Por eso al Branch and Bound se le acredita \emph{no crear cambio}
—corta el grafo— y al sorteo aleatorio \emph{no dejar patrón}.
\end{clave}

El resumen: UTXO ofrece mejor privacidad por omisión, y cuentas ofrece peor
privacidad por omisión pero mucha mejor auditabilidad, que es exactamente lo que
un contrato inteligente necesita.

% ===========================================================================
\section{Síntesis}\label{sec:sintesis}
% ===========================================================================

\begin{center}
\small
\begin{tabularx}{\textwidth}{@{}lXX@{}}
\toprule
Dimensión & UTXO & Cuentas \\
\midrule
Unidad de estado & Objeto inmutable con dueño & Celda mutable indexada por dirección \\
Operación & Destruir y crear & Sobrescribir \\
Antirrepetición & Implícita: el objeto desaparece & Explícita: contador nonce \\
Orden entre pagos & Libre si las monedas son distintas & Total y estricto por cuenta \\
Paralelismo & Alto por construcción & Requiere ejecución optimista \\
Costo de transferencia & Variable, depende de la fragmentación & Fijo, 21\,000 de gas \\
Privacidad por omisión & Mejor, erosionada por heurísticas & Peor, dirección persistente \\
Estado con contratos & Incómodo & Natural \\
Complejidad de cartera & Alta: selección de monedas & Baja: leer un saldo \\
Crecimiento del estado & Salidas de polvo permanentes & Cuentas y almacenamiento \\
\bottomrule
\end{tabularx}
\end{center}

\subsection{Más allá de la dicotomía}

Presentar esto como una elección binaria era razonable en 2018. Ya no lo es.

\textbf{Cardano} usa eUTXO: conserva los objetos, pero cada salida puede llevar
un dato arbitrario adjunto, lo que permite contratos manteniendo la validación
local. El precio es que el estado compartido se vuelve un cuello de botella.

\textbf{Sui} generaliza el UTXO a objetos con versión: los objetos propios van
por una vía rápida sin consenso completo; los compartidos sí pasan por consenso.

\textbf{Aptos y Monad} mantienen cuentas y atacan el problema desde la ejecución:
corren todo en paralelo suponiendo que no habrá conflictos y reejecutan los que
sí ocurren.

\begin{clave}{La moraleja}
Lo que importa no es la etiqueta, sino qué tan fina es la \textbf{granularidad
del estado} y \textbf{quién puede tocarlo}. UTXO y cuentas son los dos extremos
de ese eje.
\end{clave}

% ===========================================================================
\section{Ejercicios}\label{sec:ej}
% ===========================================================================

\begin{enumerate}[leftmargin=1.8em,itemsep=0.8em]

\item\label{e:enum} \textbf{Enumerar a mano.}
Una cartera tiene monedas de 1, 3, 7, 12 y 20 mil satoshis. Se quiere pagar 15
mil a 20 sat/vB. Enumera los subconjuntos válidos, calcula el desperdicio de cada
uno y decide cuál elegiría Bitcoin Core.

\item\label{e:orden} \textbf{Independencia del orden.}
Demuestra que si dos transacciones UTXO tienen conjuntos de entradas disjuntos,
el resultado no depende del orden. Muestra con un contraejemplo que la propiedad
análoga es falsa en el modelo de cuentas.

\item\label{e:grafo} \textbf{El grafo de transacciones.}
Modela el conjunto UTXO como un grafo dirigido donde cada transacción es un
vértice. ¿Qué tipo de grafo resulta? ¿Puede tener ciclos? Justifica usando que el
identificador de una transacción es el hash de su contenido.

\item\label{e:replay} \textbf{Un mundo sin nonce.}
Si el nonce no existiera en Ethereum, diseña un ataque concreto de repetición.
Después explica por qué el mismo ataque es imposible en Bitcoin sin invocar
ninguna propiedad criptográfica adicional.

\item\label{e:pruebas} \textbf{Tamaño de las pruebas.}
Una prueba de acceso en el árbol hexario ronda 3 kB. Estima el tamaño con un
árbol binario de la misma profundidad efectiva y explica de dónde sale el factor
de reducción que reporta el EIP-7864.

\item\label{e:crec} \textbf{Discusión.}
¿El crecimiento irreversible del conjunto UTXO es un problema peor o mejor que el
crecimiento del estado de contratos en Ethereum? Considera quién paga el costo.

\end{enumerate}

% ===========================================================================
\appendix
\section{Soluciones}\label{app:sol}
% ===========================================================================

\subsection*{Ejercicio \ref{e:enum}}
Hay 20 subconjuntos que alcanzan. El costo del cambio a 20 sat/vB es
$31 \times 20 + 68 \times 10 = 1300$ sat. Los tres mejores aparecen en la tabla
de la sección~\ref{sec:seleccion}: gana $\{7000, 12\,000\}$ sin cambio, con
desperdicio 1800.

La opción de una sola entrada paga menos comisión hoy (2820 frente a 4000) y aun
así pierde, porque arrastra el costo del cambio.

\subsection*{Ejercicio \ref{e:orden}}
\textbf{UTXO.} Sean $t_1, t_2$ con entradas $E_1 \cap E_2 = \emptyset$ y salidas
$O_1, O_2$ distintas entre sí (el outpoint incluye el txid). Aplicarlas en
cualquier orden da
\[
  \bigl(U \setminus (E_1 \cup E_2)\bigr) \cup O_1 \cup O_2,
\]
porque la diferencia y la unión de conjuntos conmutan cuando los conjuntos
retirados son disjuntos. La validez también se conserva: $E_2 \subseteq U$ sigue
cumpliéndose tras quitar $E_1$, justo porque son disjuntos.

\textbf{Cuentas.} Alicia tiene saldo 5. $t_1$ paga 4 a Beto, $t_2$ paga 3 a
Caro. Por separado cada una es válida. En el orden $(t_1,t_2)$ la segunda falla
por saldo insuficiente; en el inverso falla la otra. El estado final difiere. Y
eso antes de considerar el nonce, que directamente prohíbe el orden inverso.

\subsection*{Ejercicio \ref{e:grafo}}
Resulta un \textbf{grafo acíclico dirigido} (DAG). Hay una arista de $t_1$ a
$t_2$ cuando $t_2$ gasta una salida de $t_1$, y entonces $t_2$ contiene el txid
de $t_1$ entre sus entradas.

Un ciclo exigiría que el contenido de $t_1$ incluyera el hash de $t_2$, cuyo
contenido a su vez incluye el hash de $t_1$: una preimagen circular. Construirlo
equivale a romper la resistencia a preimagen de la función hash. Los UTXO
vigentes son exactamente las \emph{hojas} del DAG.

\subsection*{Ejercicio \ref{e:replay}}
\textbf{El ataque.} Alicia paga 2 ETH a una tienda. La transacción firmada es
pública. Sin nonce, la tienda —o cualquiera que lea la cadena— puede reenviarla,
byte por byte, tantas veces como quiera. Cada reenvío es válido: la firma
verifica, el origen es correcto y el saldo alcanza. Alicia pierde 2 ETH por
repetición hasta quedarse en cero. No hace falta falsificar nada; basta copiar.

\textbf{Por qué no aplica en Bitcoin.} La transacción consume outpoints
concretos. Tras el primer minado esos outpoints ya no están en $U$, así que la
copia falla en la primera regla de validez —«cada entrada existe en $U$»— antes
de que nadie mire la firma. Es un argumento puramente estructural sobre
conjuntos.

\subsection*{Ejercicio \ref{e:pruebas}}
La prueba es la lista de hermanos en cada nivel. En un árbol de aridad $k$ cada
nodo del camino aporta $k-1$ hermanos de 32 bytes, y para cubrir el mismo espacio
de claves hacen falta $\log_k N$ niveles:
\[
  (k-1)\,\log_k N \cdot 32\,\text{B}.
\]
Para $k = 16$ el factor $(k-1)/\log_2 k$ es $15/4 = 3.75$; para $k = 2$ es
$1/1 = 1$. La razón entre ambos es $\approx 3.75$, de donde sale el «unas cuatro
veces más cortas». Con 3 kB de partida quedan unos 800 bytes.

Lo que se paga a cambio es profundidad: más niveles, más hashes al verificar. Es
un intercambio de ancho de banda por cómputo, y el ancho de banda era el recurso
escaso.

\subsection*{Ejercicio \ref{e:crec}}
No tiene respuesta única. Elementos que deberían aparecer:

\begin{itemize}
  \item En UTXO el costo recae por completo en quien opera un nodo, y es una
        externalidad: quien creó el polvo ya pagó su comisión y no vuelve.
  \item En Ethereum el gas de almacenamiento internaliza parte del costo y los
        reembolsos premian liberar estado, aunque siguen siendo insuficientes.
  \item El conjunto UTXO tiene un techo natural ligado al número de salidas
        económicamente viables; el estado de contratos no lo tiene.
  \item Las soluciones van en direcciones distintas: acumuladores (Utreexo) para
        no guardar el conjunto, y caducidad de estado del lado de cuentas.
\end{itemize}

% ===========================================================================
\section{Referencias}
% ===========================================================================

\begin{itemize}[leftmargin=1.4em]
  \item Nakamoto, S. \emph{Bitcoin: A Peer-to-Peer Electronic Cash System}, 2008,
        secciones 9 y 11.
  \item Wood, G. \emph{Ethereum: A Secure Decentralised Generalised Transaction
        Ledger} (Yellow Paper), secciones 4.1 y 6.
  \item Erhardt, M. (Murch). \emph{An Evaluation of Coin Selection Strategies},
        tesis de maestría, 2016. Origen del Branch and Bound.
  \item Bitcoin Core, \texttt{src/wallet/coinselection.cpp}. Implementación de
        BnB, Knapsack, SRD y CoinGrinder.
  \item EIP-155 (repetición entre cadenas), EIP-7864 (árbol binario unificado) y
        EIP-6800 (Verkle, estancado).
  \item Chakravarty, M. et al. \emph{The Extended UTXO Model}, 2020.
  \item Gelashvili, R. et al. \emph{Block-STM}, 2022.
  \item Meiklejohn, S. et al. \emph{A Fistful of Bitcoins}, IMC 2013. La
        heurística de propiedad común de entradas.
  \item Dryja, T. \emph{Utreexo}, 2019.
  \item Fuente temática: \url{http://math.com.mx/u_vs_c.html}
\end{itemize}

\end{document}
